\section{Modelo de texto Qwen3: la naturalización del Reasoning}

\subsection{Las mejoras principales de la serie Qwen3}

En comparación con Qwen 2.5, ampliamente utilizado, las principales mejoras de Qwen3 incluyen:

\begin{enumerate}
    \item \textbf{Mejora integral de las capacidades}: una evolución integral como modelo base de nueva generación
    \item \textbf{Salto en la capacidad de Reasoning}: el cambio más notable entre 2024 y 2025
    \item \textbf{Soporte multilingüe}: compatible con 119 idiomas y dialectos
    \item \textbf{Contexto largo}: ya alcanza más de 1M de tokens; en experimentos internos se ha llegado a varios millones
\end{enumerate}

\subsection{El camino hacia la naturalización del Reasoning}

\begin{importantbox}{La unificación de Thinking e Instruct}
Lin Junyang (林俊旸) compartió una señal importante del mercado: después de lanzar en abril la versión combinada de Thinking+Instruct, \textbf{más del 90\% de los clientes abandonaron el modelo Thinking independiente} y volvieron al modelo Instruct. Las razones fueron:
\begin{itemize}
    \item Al principio a los usuarios les gustaba ver el «autodiálogo» del modelo (thinking trace), pero esa novedad se desvaneció rápido
    \item El modelo Instruct es más práctico y sus respuestas son más naturales
    \item Lo importante es cómo lograr que la capacidad de Reasoning \textbf{se integre de manera natural} al modelo Instruct
\end{itemize}
\end{importantbox}

El equipo de Qwen lanzó una versión mejorada en julio; el enfoque central consistió en integrar el proceso de Thinking directamente en la respuesta para las tareas Reasoning-intensive, en lugar de usar tokens de thinking independientes. El resultado fue que puntos de referencia matemáticos como AIME subieron de alrededor de 20 puntos a alrededor de 70 puntos (el modelo Thinking puede llegar a 90 puntos), y los clientes reportaron que el modelo «se volvió notablemente más inteligente».

\begin{knowledgebox}{Las lecciones y la reflexión sobre DeepSeek R1}
DeepSeek R1 demostró que un «autodiálogo parlanchín» puede hacer más fuerte al modelo, pero el equipo de Qwen considera que esa no es la forma final. Lo que persiguen es un \textbf{Reasoning más natural}: el modelo sí piensa cuando de verdad necesita reflexionar en profundidad, pero no interrumpe la experiencia del usuario con un monólogo interno extenso.
\end{knowledgebox}

\subsection{Multilingüismo: del coreano al urdu}

El origen del soporte multilingüe es una historia interesante:

\begin{enumerate}
    \item Al principio Qwen solo era bilingüe en chino e inglés
    \item Un investigador coreano dijo: «Su modelo no entiende nada de coreano», y eligió Mistral como modelo base
    \item Al investigar se descubrió que el modelo Pretrain en realidad sí sabía coreano, solo que el Post-train no lo cubría
    \item Lo corrigieron de paso, y empezaron a llegar usuarios de todo el mundo
    \item Usuarios de Pakistán pidieron una y otra vez soporte para el urdu: «De verdad no tenemos ningún modelo grande disponible»
    \item Actualmente ya se admiten 119 idiomas, pero la cobertura de idiomas africanos sigue siendo insuficiente
\end{enumerate}

\begin{warningbox}{La realidad de la brecha digital}
Lin Junyang mencionó que en muchas regiones de África todavía se usan teléfonos básicos. Si la intención de construir un modelo grande no es ayudar a toda la humanidad, «entonces mejor no hacerlo». Esta postura idealista impulsa el esfuerzo continuo de Qwen por ampliar la cobertura de idiomas.
\end{warningbox}

\subsection{Contexto largo y Memory}

\begin{knowledgebox}{La relación entre Context Length y RAG}
Lin Junyang citó la opinión del leader del equipo de secuencias largas de Gemini: Context Length y RAG son, en esencia, \textbf{ortogonales}.
\begin{itemize}
    \item No se debería meter una gran cantidad de «basura» en el contexto largo
    \item Un contexto más largo sienta las bases para una mejor gestión de la memoria
    \item Qwen ya alcanza más de 1M de tokens; en experimentos internos se llega a varios millones
    \item Uno de los motivos para cambiar la arquitectura del modelo (Qwen3-Next) es explorar la posibilidad de un infinite long context
\end{itemize}
\end{knowledgebox}

\subsection{Resumen de la sección}

El avance central del modelo de texto Qwen3 consiste en integrar de manera natural la capacidad de Reasoning al modelo Instruct, a la vez que se amplían el soporte multilingüe y la capacidad de contexto largo. «Hacer que el modelo sea inteligente pero no molesto» es el equilibrio clave de la conversión en producto.

